\begin{animation}[id="bfs-queue", label="BFS Queue Demo"]
\shape{q}{Queue}{capacity=6, data=[1], label="$Q$"}

\step
\recolor{q.cell[0]}{state=current}
\narrate{BFS starts. Enqueue source node 1. Front and rear both point to cell 0.}

\step
\apply{q}{enqueue=2}
\apply{q}{enqueue=4}
\recolor{q.cell[0]}{state=idle}
\recolor{q.cell[1]}{state=good}
\recolor{q.cell[2]}{state=good}
\annotate{q.cell[1]}{label="enq", arrow_from="q.cell[0]", color=good}
\annotate{q.cell[2]}{label="enq", arrow_from="q.cell[1]", color=good}
\narrate{Discover neighbors of node 1. Enqueue nodes 2 and 4.}

\step
\apply{q}{dequeue=true}
\recolor{q.cell[0]}{state=dim}
\recolor{q.cell[1]}{state=current}
\recolor{q.cell[2]}{state=idle}
\annotate{q.cell[0]}{label="deq", arrow_from="q.cell[1]", color=error}
\narrate{Dequeue node 1 (done). Process node 2 next. Front pointer advances.}

\step
\apply{q}{dequeue=true}
\apply{q}{enqueue=5}
\apply{q}{enqueue=3}
\recolor{q.cell[1]}{state=dim}
\recolor{q.cell[2]}{state=current}
\recolor{q.cell[3]}{state=good}
\recolor{q.cell[4]}{state=good}
\narrate{Dequeue node 2 (done). Process node 4. Enqueue its neighbors 5 and 3.}

\step
\apply{q}{dequeue=true}
\recolor{q.cell[2]}{state=dim}
\recolor{q.cell[3]}{state=current}
\recolor{q.cell[4]}{state=idle}
\narrate{Dequeue node 4 (done). Process node 5 next. Front advances to cell 3.}

\step
\recolor{q.cell[0]}{state=dim}
\recolor{q.cell[1]}{state=dim}
\recolor{q.cell[2]}{state=dim}
\recolor{q.cell[3]}{state=done}
\recolor{q.cell[4]}{state=done}
\narrate{BFS complete. All reachable nodes have been visited.}
\end{animation}
